\subsection{LINGUAGENS DE PROGRAMAÇÃO}

A linguagem nativa de computadores é a binária (uns e zeros), todas as instruções e dados devem ser providas nesse formato para serem processados. Código binário nativo é chamado de linguagem de máquina. Com o passar do tempo foram aparecendo novas formas de construir programas; sucedendo o código de máquina os programas começaram a serem escritos em sequências de números hexadecimais. Em seguida o surgimento linguagens \en{assembly}, os códigos então começaram a ser escritos usando uma série de comandos mnemônicos que represetam uma conjunto de instruções em código de máquina. Essas linguagens de programação não são apropriadas para o desenvolvimento de aplicação em larga escala, portanto, o surgimento das linguagens de alto nível como C. Estas definem um conjunto de mecanismos que servem como abstrações das instruções de máquina (\en{loops}, condicionais, funções, etc) \cite{Bailey2005IntroductionC}.

  Construir estruturas de dados e algoritimos requer uma linguagem de programação de alto nível \cite{GoodrichTamassiaMount2011DataStructuresCpp}. Logo, para o desenvolvimento da aplicação móvel e dos \en{middlewares} de sincronização foram, portanto, escolhidas as linguagens de programação \en{Kotlin} e \en{Golang}, respectivamente.  
